\section{Introdução}
\label{sec:introducao}

% PRISMA Item 3: Rationale — justificativa no contexto do conhecimento existente
% PRISMA Item 4: Objectives — declaração explícita das questões abordadas

Nove dos 28 estudos sobre persistência de conhecimento em
agentes de codificação relatam cenários em que a memória
prejudica o desempenho do agente.
Esse achado torna crítico distinguir quando persistir conhecimento
é vantajoso e quando é prejudicial, decisão que precede a
comparação entre arquiteturas tecnicamente sofisticadas.
Agentes de codificação baseados em modelos de linguagem de grande
escala (LLMs) alcançaram adoção crescente em tarefas de Engenharia
de Software como resolução automatizada de issues, geração de
código e manutenção de repositórios~\cite{du2026memory}.
Sistemas como SWE-Agent, OpenHands e Agentless interagem
iterativamente com o código-fonte, executam testes e propõem
correções sem intervenção humana direta.
A capacidade desses agentes, contudo, permanece limitada por uma
restrição básica: cada sessão parte do zero.

O agente redescobre a arquitetura do projeto, repete erros já
cometidos em sessões anteriores e ignora decisões de design
previamente validadas.
Essa volatilidade de conhecimento equivale a um desenvolvedor sem
memória entre reuniões.
A persistência de conhecimento entre sessões surgiu como resposta
a essa limitação: ao registrar informações de resoluções
anteriores, o agente pode reutilizá-las nas tarefas seguintes.

Ao menos 24 sistemas de persistência foram propostos desde 2024,
todos publicados entre 2024 e o primeiro trimestre de 2026.
As arquiteturas variam de bancos de experiência que destilam
trajetórias de resolução a grafos de conhecimento, memória guiada
por AST e metáforas de controle de versão inspiradas no Git.
O campo, porém, permanece fragmentado: os estudos não convergem
para uma taxonomia unificada, e a maioria não cita trabalhos
contemporâneos sobre o mesmo problema.

Dessa fragmentação emergem lacunas que limitam tanto a pesquisa
quanto a prática.
Nenhum estudo compara três ou mais arquiteturas de persistência no
mesmo benchmark com o mesmo modelo; toda evidência de superioridade
arquitetural depende de comparações entre estudos com condições
experimentais heterogêneas.
A subnotificação de custos agrava o problema: 43\% dos estudos
omitem dados de custo, e nenhum constrói uma fronteira de Pareto
entre acurácia e custo computacional.
Tampouco existe evidência controlada sobre a relação entre
complexidade arquitetural e eficácia, de modo que não se sabe se
arquiteturas simples são suficientes ou se a sofisticação se
justifica em cenários específicos.

Trabalhos secundários existentes cobrem aspectos adjacentes sem
preencher essas lacunas; o posicionamento detalhado desta revisão
diante deles é apresentado na
\autoref{sec:trabalhos-relacionados}.

Esta revisão sistemática responde à seguinte questão principal: \textit{de que forma a literatura existente aborda a
persistência de conhecimento em agentes de codificação baseados em
IA, e que evidências empíricas existem sobre a eficácia e
custo-eficiência de diferentes arquiteturas de persistência?}
Para operacionalizar essa questão principal, cinco questões
secundárias orientam a síntese, cada uma motivada por uma lacuna
específica no estado atual do campo.

\noindent\textbf{QS1.} Quais arquiteturas de persistência foram
propostas, e como se classificam quanto a escopo temporal,
substrato representacional e política de controle?
A heterogeneidade arquitetural impede tanto a comparação direta
quanto a adoção informada por equipes que precisam escolher entre
alternativas concorrentes.
Uma taxonomia explícita revela quais escolhas de projeto dominam
o campo e quais permanecem inexploradas.

\noindent\textbf{QS2.} Quais métodos de avaliação e benchmarks são
utilizados, e quão comparáveis são os resultados?
Quando estudos diferentes adotam benchmarks, modelos base e
protocolos experimentais distintos, ganhos reportados não são
diretamente comparáveis entre si.
Esta questão verifica se há convergência metodológica e identifica
os obstáculos que impedem a síntese cross-study.

\noindent\textbf{QS3.} Quais resultados de desempenho são
reportados em relação a baselines sem persistência?
A magnitude dos ganhos determina se persistência representa uma
melhoria marginal ou substantiva sobre os agentes sem memória.
Sintetizar esses resultados expõe se o ganho se mantém em
diferentes condições experimentais ou depende de contextos
específicos.

\noindent\textbf{QS4.} Quais métricas de custo e eficiência são
reportadas, e qual a relação entre custo e eficácia?
Persistência adiciona overhead de armazenamento, recuperação e
contexto efetivo do agente.
Sem dados de custo, qualquer ganho de acurácia é incompleto para
fins de adoção prática, e a relação custo-eficácia permanece
inacessível ao leitor.

\noindent\textbf{QS5.} Que evidências existem sobre a relação
entre complexidade arquitetural e eficácia, incluindo degradação?
Sistemas mais sofisticados nem sempre superam alternativas
simples, e mais memória nem sempre produz mais eficácia.
Identificar mecanismos de degradação e padrões de complexidade
versus ganho informa decisões de projeto em cenários onde a
sofisticação não se justifica.

\begin{table}[htbp]
\centering
\caption{Mapa das questões secundárias: objetivo de cada questão e
localização da resposta no manuscrito.}
\label{tab:mapa-rqs}
\small
\begin{tabularx}{\textwidth}{l X l l}
\toprule
\textbf{QS} & \textbf{Objetivo} & \textbf{Seção} & \textbf{Tabela/Figura} \\
\midrule
QS1 & Mapear e classificar as arquiteturas de persistência &
  \autoref{sec:qs1} & \autoref{tab:arquiteturas}, \autoref{fig:taxonomia} \\
QS2 & Catalogar métodos de avaliação e benchmarks &
  \autoref{sec:qs2} & \autoref{tab:benchmarks} \\
QS3 & Sintetizar resultados de desempenho vs.\ baselines &
  \autoref{sec:qs3} & \autoref{tab:desempenho}, \autoref{fig:ganhos} \\
QS4 & Quantificar cobertura e relação custo-eficácia &
  \autoref{sec:qs4} & \autoref{tab:custo}, \autoref{fig:custo} \\
QS5 & Avaliar relação entre complexidade e eficácia &
  \autoref{sec:qs5} & \autoref{tab:degradacao} \\
\bottomrule
\end{tabularx}
\end{table}

\noindent O escopo abrange publicações de 2023 a março de 2026, com
buscas em Scopus e arXiv complementadas por rastreamento de citações.

As contribuições desta revisão são as seguintes:

\begin{enumerate}
  \item o primeiro mapeamento sistemático do campo, classificando
    24 arquiteturas de persistência para agentes de codificação
    segundo três dimensões: escopo temporal da memória, substrato
    representacional e política de leitura e escrita;
  \item a síntese quantitativa de 19 comparações controladas sobre
    eficácia e custo, com identificação de rendimentos decrescentes
    em baselines fortes e de impacto bimodal sobre o custo
    operacional;
  \item a confirmação empírica de três lacunas centrais do campo:
    ausência de comparações head-to-head entre arquiteturas no
    mesmo benchmark, subnotificação sistemática de custos por
    tarefa e cobertura insuficiente da relação entre complexidade
    arquitetural e eficácia.
\end{enumerate}

As lacunas confirmadas motivam diretamente o estudo comparativo
empírico planejado como segunda etapa desta dissertação, descrito
em mais detalhe na \autoref{sec:direcoes-futuras}.
